\section{November 2026: the newest report and the examiner Hoedt}

\begin{summary}{What was asked, and what it means for the committee Zoitl, Hoedt, Pirker}
\textsf{November 2026} (chair Kofler, ML\&P Küng, AI\&LS Luukkonen): the student wrote the maximum-likelihood
formula instead of the empirical risk and was guided to it. Then: what happens when $N\to\infty$, and what makes
deep networks powerful from a computational point of view (``nested functions''). The defense was strictly
timed. \
\textsf{September 2026} (ML\&P Hoedt): what is a neural network and a single neuron, why ``network'', why a
non-linear activation, why deep networks instead of simpler models, then backprop and $\delta^{[l]}$. \
\textsf{Rule:} write the simplest correct formula first and let the examiner steer. Answer from the examiner's
own research: Hoedt wrote the chapters on learning methods, initialization and higher-order optimization of the
DLNN script. Full answers with audio: \url{https://eigenzayn.github.io/committee-exam-prep/jawaban.html}.
\end{summary}

\begin{qabox}[Küng]{\asked Write down empirical risk minimization. How is it different from maximum likelihood?}
\from{DL Basic Techniques (DLNN I) · Ch.~5 ERM, Ch.~4 linear model · asked November 2026}
\ans ERM picks the parameters that minimise the average loss on the training data. Maximum likelihood picks
the parameters that make the observed labels most probable. They coincide when the loss is the negative log
of the model's probability, for example squared loss with Gaussian noise or cross-entropy with a Bernoulli
or softmax output.
\formula Step 1, the empirical risk, for any model $g$ and any loss $L$:
\[ \Remp(\mathbf w)=\frac1N\sum_{n=1}^N L\bigl(y_n,g(\mathbf x_n;\mathbf w)\bigr),\qquad \hat{\mathbf w}=\arg\min_{\mathbf w}\Remp(\mathbf w). \]
Step 2, the likelihood of independent labels, and its negative logarithm:
\[ p(\mathbf y\mid X,\mathbf w)=\prod_{n=1}^N p(y_n\mid\mathbf x_n,\mathbf w),\qquad
-\log p=\sum_{n=1}^N -\log p(y_n\mid\mathbf x_n,\mathbf w). \]
Step 3, choose $L(y,g)=-\log p(y\mid\mathbf x,\mathbf w)$. Then $-\log p = N\,\Remp$, so maximising the likelihood
is minimising the empirical risk. Gaussian noise gives $-\log p=\frac N2\log(2\pi\sigma^2)+\frac1{2\sigma^2}\sum_n(y_n-g_n)^2$:
the squared loss.
\follow ``Is every loss a likelihood?'' No. The zero-one loss or a hinge loss is not the negative log of a
probability model, but it is still a valid loss for ERM.
\end{qabox}

\begin{qabox}[Küng]{\asked What happens when the number of samples goes to infinity?}
\from{DL Basic Techniques (DLNN I) · Ch.~5 generalization error, test set · asked November 2026}
\ans For a fixed model the empirical risk converges to the true risk, by the law of large numbers. The
catch is that the parameters were chosen on the same data, so the training error is optimistic; that is
why we need a separate test set.
\formula Risk and its estimate:
\[ \Rgen(\mathbf w)=\mathbb E_{(\mathbf x,y)\sim p}\bigl[L(y,g(\mathbf x;\mathbf w))\bigr]
=\int L(y,g(\mathbf x;\mathbf w))\,p(\mathbf x,y)\,\mathrm d\mathbf x\,\mathrm dy . \]
For fixed $\mathbf w$ the losses $L_n$ are i.i.d.\ with mean $\Rgen(\mathbf w)$, so
\[ \Remp(\mathbf w)=\frac1N\sum_n L_n\;\xrightarrow{N\to\infty}\;\Rgen(\mathbf w),\qquad
\operatorname{Var}\bigl(\Remp\bigr)=\frac{\operatorname{Var}(L)}{N}. \]
The error of the estimate shrinks like $1/\sqrt N$.
\follow ``And for $\hat{\mathbf w}$?'' With more data the gap between training and test error closes,
because a fixed model class can fit less of the noise.
\end{qabox}

\begin{qabox}[Küng]{\asked What makes deep networks so powerful from a computational perspective?}
\from{DL Basic Techniques (DLNN I) · Ch.~4 deep feed-forward networks, Eq.~(4.74) · asked November 2026}
\ans A deep network is a composition of simple functions, a function of a function of a function. Each
layer reuses what the layer below computed, and the chain rule over this composition gives all gradients in
one backward pass, at a cost linear in the number of weights.
\formula The nested function and the cost:
\[ \hat{\mathbf y}=\sigma\Bigl(W^{[L]}\,f\bigl(\cdots f\bigl(W^{[2]}f(W^{[1]}\mathbf x)\bigr)\cdots\bigr)\Bigr),\qquad
\text{forward and backward pass: } O(W). \]
Backward along the composition: $\boldsymbol\delta^{[l-1]T}=\boldsymbol\delta^{[l]T}W^{[l]}\operatorname{diag}\bigl(f'(\mathbf s^{[l-1]})\bigr)$, and
$\partial L/\partial W^{[l]}$ uses $\boldsymbol\delta^{[l]}$ and $\mathbf a^{[l-1]}$.
\follow ``Is one hidden layer not enough?'' In principle yes, by the universal approximation theorem, but the
theorem does not say how many units are needed, and composition reuses features instead.
\end{qabox}

\begin{qabox}[Hoedt]{\asked What is a neural network, what is a single neuron, and why is it a ``network''?}
\from{DL Basic Techniques (DLNN I), lecturers Klambauer and Hoedt · Ch.~4 · asked September 2026}
\ans A neuron forms a weighted sum of its inputs plus a bias and passes it through an activation function. A
network connects many neurons so that the output of one becomes the input of others; in a feed-forward net they
are arranged in layers.
\formula \[ s_i=\sum_j w_{ij}a_j+b_i,\qquad a_i=f(s_i),\qquad \mathbf a^{[l]}=f\bigl(W^{[l]}\mathbf a^{[l-1]}+\mathbf b^{[l]}\bigr). \]
\end{qabox}

\begin{qabox}[Hoedt]{\asked Why is a non-linear activation necessary, and why deep networks instead of simpler models?}
\from{DL Basic Techniques (DLNN I) · Ch.~4 XOR; DLNN II Ch.~13 universal approximation · asked September 2026}
\ans Without a non-linearity the layers collapse into one linear map, which cannot even solve XOR. With
one, a network can approximate any continuous function on a compact set, it reuses features through
composition, and backpropagation makes it trainable at scale.
\formula Collapse: $W_2(W_1\mathbf x+\mathbf b_1)+\mathbf b_2=W'\mathbf x+\mathbf b'$. XOR with $g=x_1w_1+x_2w_2$: inputs $(1,0)$ and $(0,1)$ give
$w_1=w_2=1$, then $(1,1)$ demands $0=2$. With ReLU, $W=\begin{pmatrix}1&1\\1&1\end{pmatrix}$, $\mathbf b=\begin{pmatrix}0\\-1\end{pmatrix}$,
$\mathbf w=\begin{pmatrix}1\\-2\end{pmatrix}$ maps $(0,0),(1,0),(0,1),(1,1)$ to $0,1,1,0$.
\follow Then define $\delta^{[l]}=\partial L/\partial\mathbf s^{[l]}$ and derive the recursion (Chapter on backprop).
\end{qabox}

\begin{qabox}[Defense]{\asked Advice from the November 2026 defense}
\from{report of November 2026 (chair Kofler)}
\ans Keep the presentation inside the time, because the committee was strict and the student had to skip
slides. The committee asks about what it does not understand, since you are the expert on your topic. A
typical question is what you would do differently if you started again.
\follow For this thesis: ``I would plan a second mesh and a comparison with measurements from the start,
and record only the kinematic part of the flux so that it can transfer.''
\end{qabox}
